%==============================================================================
\chapter{Host (DAW) compatibility}\label{ch:daw}
%==============================================================================

SuperMoTo is primarily designed and tested in \textbf{REAPER}. To host it usefully a DAW must satisfy two requirements:

\begin{enumerate}[nosep]
  \item \textbf{Flexible channel routing}: the ability to feed the plugin's discrete inputs from arbitrary sources and to send its discrete outputs to arbitrary hardware outputs / busses.
  \item \textbf{A monitoring (control-room) insert point}: a place to put the plugin on the path between the mix and the physical speaker outputs, after the program material is mixed but before it leaves for the monitors, ideally not printed into renders/bounces.
\end{enumerate}

\section{Recommended and likely-compatible hosts}
\begin{table}[ht]
\centering
\renewcommand{\arraystretch}{1.3}
\begin{tabularx}{\textwidth}{l c X}
\toprule
\textbf{Host} & \textbf{Fit} & \textbf{Notes} \\
\midrule
REAPER & \textbf{Primary} &
  Per-track/track-channel routing up to 64 channels, and a dedicated
  \emph{Monitoring FX} chain (\textsf{View\,$\rightarrow$\,Monitoring FX})
  that sits on the hardware output path and is excluded from renders. The
  reference environment. \\
Ardour & Strong &
  Arbitrary bus channel counts and flexible port routing; a dedicated
  \emph{Monitor Section} on the master/monitor path. A natural fit on Linux. \\
Steinberg Nuendo / Cubase & Strong &
  The \emph{Control Room} provides a monitor insert chain with flexible
  downmix and multi-channel busses, ideal for a post-production monitoring
  matrix. Confirm the Control Room insert accepts a 32-channel plugin layout. \\
Bitwig Studio & Possible &
  Very flexible routing and per-device chains; place the plugin on the master
  or on a hardware-output chain. No formal ``monitor section'', so the
  monitoring insert must be built manually and bypassed before export. \\
Standalone (built in) & Always &
  The Standalone build talks straight to the audio device (up to 32 in / 32 out
  if the interface provides them); useful for a dedicated monitor-controller
  machine with no DAW. \\
\bottomrule
\end{tabularx}
\caption{Hosts that can drive SuperMoTo. "Fit" is the author's assessment
based on the routing / monitor-section requirements; only REAPER is formally
tested.}
\label{tab:hosts}
\end{table}

\section{Hosts likely to be awkward}
\begin{itemize}[nosep]
  \item \textbf{Pro Tools}: plugin channel formats are constrained to named
        surround layouts; an arbitrary high-count discrete insert is not the
        native model.
  \item \textbf{Logic Pro}: surround formats are fixed and discrete
        32-channel routing through a monitor insert is not straightforward.
  \item \textbf{Studio One / FL Studio / Live}: usable as generic effect
        hosts, but they lack a true control-room monitor insert, so the
        plugin would have to live on the master bus and be manually bypassed
        for exports.
\end{itemize}

\section{Suggested REAPER setup (reference)}
\begin{enumerate}[nosep]
  \item Open \textsf{View\,$\rightarrow$\,Monitoring FX} and insert SuperMoTo
        there so it processes the hardware output path and is excluded from
        renders.
  \item In the plugin's pin connector / channel mapping, map your mix busses
        to inputs 1..N and the plugin outputs 1..M to the physical speaker
        outputs.
  \item Build your routing in the matrix (or with the Configuration tool) and
        store rigs in configurations A..F.
\end{enumerate}

% TODO(screenshot): REAPER Monitoring FX window with SuperMoTo inserted and
% the 32x32 pin connector visible.
\screenshot{reaper-monitoring-fx.png}{REAPER ``Monitoring FX'' window with
SuperMoTo inserted, showing the fixed 32-in/32-out pin/channel mapping to the
hardware outputs.}
